\section{Background}

\subsection*{Trajectory methods systematically miss cyclic cell-state structure}

Single-cell RNA sequencing has transformed our view of cellular
heterogeneity, but the dominant trajectory toolkit for ordering cells
--- Monocle and Slingshot (tree backbones), PAGA (connectivity graph),
and RNA velocity (vector field) --- was designed to summarise
progression rather than to quantify reversible cycling
\citep{saelens2019}. PAGA admits non-tree graphs and velocity defines
a field, but neither yields a scalar that scores how robustly cells
trace a closed loop across analysis scales. That gap matters whenever
cell populations reversibly transition between states --- as in the
epithelial-to-mesenchymal transition (EMT), the cell cycle, and
adaptive drug tolerance \citep{shaffer2017,oren2021} --- because a
reversible cycle is mathematically a one-dimensional hole in the
data's point cloud.

\subsection*{Topological data analysis as a framework-level remedy}

\paragraph{An intuitive picture for biomedical readers.}
Drug-tolerant cancer cells appear not to die or jump to one resistant
state but to reversibly cycle between several transcriptional states
(\cref{fig:concept}a). When each cell is plotted as a point in
gene-expression space, this reversible cycling traces a closed loop
--- the same kind of one-dimensional ``hole'' that distinguishes a
donut from a sphere in topology, and which tree-based trajectory
methods discard by force-fitting branching graphs
(\cref{fig:concept}b). Persistent homology quantifies the loop's
robustness across spatial scales: the lifespan of the longest such
loop is the \emph{maximum $H_1$ persistence} statistic used
throughout this paper. The headline finding is that this statistic
rises monotonically with EGFR-TKI exposure across PC9 cell line,
two PDX models, and a 14-patient longitudinal cohort
(\cref{fig:concept}c). Mapper complements this with an interpretable
graph in which nodes group transcriptionally similar cells and edges
connect related nodes \citep{singh2007}.

Formally, persistent homology returns a multiset of birth--death
pairs (the persistence diagram) over a filtered sequence of simplicial
complexes \citep{edelsbrunner2010}, and Mapper \citep{singh2007}
returns a complementary graph summary. Both have been applied
individually to bulk cancer transcriptomes \citep{nicolau2011} and
developmental scRNA-seq \citep{rizvi2017}; prior work has not
integrated them into a reproducible, open-source framework for
longitudinal scRNA-seq with formal null-model testing and standard
controls --- the gap our contribution addresses.

\subsection*{Drug-tolerant persisters as a testbed}

Drug-tolerant persister cells (DTPs) in EGFR-mutant non-small cell
lung cancer (NSCLC) provide an ideal testbed for a topology-based
framework. First-line osimertinib gives a median progression-free
survival of roughly 18 months, yet essentially every patient relapses
from residual disease, and no approved drug specifically targets the
persister state \citep{ramalingam2020,piotrowska2018}. DTPs are a
small, slow-growing subpopulation that survives targeted therapy not
by acquiring mutations but by entering a reversible, quiescent
transcriptional state --- behaving like cells hibernating through
drug pressure and waking up when conditions relax
\citep{sharma2010,hata2016,aissa2021,oren2021}. Multiple public
scRNA-seq datasets cover longitudinal treatment with erlotinib (1st-
generation TKI) or osimertinib (3rd-generation TKI) at cell-line and
patient-derived xenograft (PDX) scales, and lineage tracing
\citep{oren2021} has established that DTPs contain cycling
subpopulations --- but the topological signature of this dynamic has
not been formally quantified.

\subsection*{Contribution}

We present an open-source pipeline that applies persistent homology
($H_0$, $H_1$) and the Mapper algorithm to longitudinal scRNA-seq,
with integrated permutation-null testing, two independent cell-cycle
controls, Wasserstein diagram comparison, and Mapper-based gene
attribution. We validate it on five independent EGFR-mutant
lung-cancer datasets covering two TKIs (erlotinib, osimertinib),
three biological scales (cell line, PDX, patient tumour), and a
14-patient longitudinal cohort. The maximum $H_1$ persistence rises
monotonically with drug exposure across the validation cohorts,
survives both cell-cycle controls, and is biologically anchored by
EMT-program enrichment in loop-associated transcripts (cell-line and
PDX overlap $p = 1.25 \times 10^{-11}$). All analyses are reproducible
from raw GEO data via a single Makefile target under MIT license.
